%%%%%%%%%%%%%%%%%%%%%%%%%%%%%%%%%%%%%%%%%%%%%%%%%%%%%%%%%%%%%%%%%%%%%%%%%%%%%%%%
% BLOQUE 4 — MARCO TEÓRICO Y CONCEPTUAL
% Sección agregada tras cruzar la ficha TI-05 con las pautas generales
% del curso, que exigen esta sección entre la Introducción y el
% Desarrollo. Da las bases conceptuales que el Desarrollo (bloque 5)
% profundiza técnicamente.
%%%%%%%%%%%%%%%%%%%%%%%%%%%%%%%%%%%%%%%%%%%%%%%%%%%%%%%%%%%%%%%%%%%%%%%%%%%%%%%%

\section{Marco teórico y conceptual}

% TODO (equipo): definiciones base necesarias para entender el resto
% del informe:
%   - Serverless (arquitectura sin servidor): qué es y qué no es.
%   - FaaS (funciones como servicio).
%   - Contenedores sin servidor.
%   - Edge computing (cómputo en el borde).
%   - El continuo nube-borde (cloud-to-edge continuum): dónde se ubica
%     cada modelo de ejecución entre el centro de datos y el borde de
%     la red.
% Usar fuentes preferentes (documentación oficial, estándares, papers
% revisados por pares) según el punto 5 de las indicaciones del curso.

\clearpage
